\documentclass[12pt, a4paper]{article}
\usepackage[utf8]{inputenc}
\usepackage[english]{babel}
\usepackage[margin=2cm]{geometry}
\usepackage{graphicx}
\usepackage{float} %use the option [H]
\graphicspath{ {images/} }
\usepackage{amsthm} %lets us use \begin{proof}
\usepackage{amsmath}
\DeclareMathOperator{\arcsinh}{arcsinh}
\usepackage{amssymb} %gives us the character
\usepackage{CJKutf8}
\usepackage[export]{adjustbox}
\setlength{\parindent}{0cm} % starting spaces
\setlength{\parskip}{1em} % paragraph width
\usepackage{array}
\usepackage{tabularx}
\usepackage{markdown}
\title{\textbf{Calculus HW2}}
\author{王弘禹}
\date\today
\begin{document}
\begin{CJK*}{UTF8}{bsmi}
\linespread{1.5}
\maketitle
\newcommand{\st}[1]{\section*{#1}}
\newcommand{\sst}[1]{\subsection*{#1}}
\newcommand{\ssst}[1]{\subsubsection*{#1}}
\newcommand{\dsp}{\displaystyle}
\newcommand{\dx}{\,dx}
\newcommand{\dphi}{\,d\phi}
\newcommand{\dtheta}{\,d\theta}
\newcommand{\dy}{\,dy}
\newcommand{\dz}{\,dz}
\newcommand{\dr}{\,dr}
\newcommand{\drho}{\,d\rho}
\newcommand{\dt}{\,dt}
\st{16.4}
\sst{22.}
\begin{align}
    \textbf{F}(x,y)=\langle\sin x,\sin y+xy^2+\dfrac{1}{3}x^3\rangle
\end{align}
By green theorem and polar coordinate:
\begin{align}
    \int_C \textbf{F}\cdot \,d\textbf{r} = \int_0^5\int_0^{\pi/2}r^3\dtheta dr=\dfrac{\pi}{2}\dfrac{625}{4}
\end{align}
\sst{31.}
\begin{equation}
    \dfrac{\partial}{\partial y}F_x=\dfrac{\partial }{\partial x}F_y = \dfrac{2x^2-6xy^2}{x^2+y^2}
\end{equation}
Therefore, by the method of example 5, we know that all the path which is closed to origin have the same integral. So, we chose a unit circle $C_1$ and use polar coordinate:
\begin{align}
    &\textbf{F} = \langle 2\sin t\cos t, \sin^2 t-\cos^2 t\rangle\\
    &\int_{C_1}\textbf{F}\cdot d\textbf{r}=\int_0^{2\pi} (-2\sin^2 t\cos t+\sin^2 t\cos t-\cos^3 t)\dt=\int_0^{2\pi}-\cos t\dt =0=\int_{C}\textbf{F}\cdot d\textbf{r}
\end{align}

\st{16.5}
\sst{20.}
The domain $D\in\textbf{R}$ is star-shaped, therefore, we just need to check these conditions
\begin{align}
    \mathop{\forall}\limits_{i\ne j} i,j\implies \dfrac{\partial F_i}{\partial j}=\dfrac{\partial F_j}{\partial i}
\end{align}
in this case
\begin{align}
    &\dfrac{\partial F_y}{\partial x} - \dfrac{\partial F_x}{\partial y} = 0-0=0\\
    &\dfrac{\partial F_x}{\partial z} - \dfrac{\partial F_z}{\partial x} = e^z\cos x - e^z \cos x=0\\
    &\dfrac{\partial F_z}{\partial y} - \dfrac{\partial F_y}{\partial z} = -e^y\sin z - (-e^y\sin z) = 0
\end{align}


\sst{29.}
\sst{34.}
\sst{35.}

\end{CJK*}
\end{document}